\chapter{Fundamental Concepts of Quantum Mechanics}

\section{The Uncertainty Principle}

\SourcePage{10}{13}
When classical mechanics and electrodynamics are applied in an attempt to
explain atomic phenomena, they give results sharply at variance with
experiment. This is already seen most clearly in the contradiction produced
by applying ordinary electrodynamics to a model of the atom in which electrons
move around the nucleus in classical orbits. Since any accelerated charge radiates, electrons
in such orbits would be obliged to emit electromagnetic waves without
interruption. Through radiation they would lose energy
and eventually fall into the nucleus. Classical electrodynamics would therefore render the atom
unstable---a conclusion entirely at odds with reality.

So profound a conflict between theory and experiment shows that a theory
applicable to atomic phenomena---phenomena involving particles of very small
mass in very small regions of space---requires a fundamental alteration of
the basic classical ideas and laws.

A convenient point of departure for identifying these changes is the
experimentally observed phenomenon known as electron diffraction\footnote{In
fact, electron diffraction was discovered after quantum mechanics had been
created. Our presentation, however, does not follow the historical sequence
in which the theory developed; it is arranged instead to show as clearly as
possible how the fundamental principles of quantum mechanics are connected
with observed phenomena.}. When a uniform beam of electrons passes through a
crystal, the transmitted beam displays a succession of intensity maxima and
minima entirely analogous to the diffraction pattern observed for
electromagnetic waves. Under certain conditions, therefore, the behavior of
material particles---electrons---shows features characteristic of wave
processes.

\SourcePage{11}{14}
The depth of this phenomenon's conflict with ordinary ideas of motion is best
shown by the following thought experiment, an idealization of electron
diffraction by a crystal. Imagine a barrier impermeable to electrons, with two
slits cut in it. If we observe an electron beam\footnote{The beam is assumed
to be so dilute that interactions among its particles play no part.} passing
through one slit while the other is closed, we obtain some intensity
distribution on a solid screen behind the slit. Opening the second slit and
closing the first gives another distribution. With both slits open at once,
ordinary reasoning would lead us to expect simply the superposition of those
two distributions: each electron, moving on its own trajectory, passes
through one slit without affecting electrons passing through the other.
Electron diffraction shows instead that the actual result is a diffraction
pattern which, because of interference, is by no means the sum of the patterns
from the separate slits. Plainly, this result cannot be reconciled in any way
with the idea that electrons move along trajectories.

The mechanics governing atomic phenomena---the so-called \emph{quantum} or
\emph{wave mechanics}---must therefore rest on concepts of motion fundamentally
different from those of classical mechanics. Quantum mechanics has no concept
of a particle trajectory. This fact constitutes the content of the so-called
\emph{uncertainty principle}, a fundamental principle of quantum
mechanics, due to \emph{Heisenberg} (\emph{W.~Heisenberg},
1927)\footnote{It is noteworthy that the complete mathematical apparatus of
quantum mechanics had already been constructed by \emph{W.~Heisenberg} and
\emph{E.~Schrödinger} in 1925--1926---prior to formulation of the uncertainty
principle, which lays bare the physical content of that apparatus.}. Since it
rejects the customary ideas of classical mechanics, the uncertainty principle
may be said to have negative content. By itself it is naturally quite
insufficient as the basis for constructing a new mechanics of particles. Such
a theory must of course rest on certain positive propositions, which will be
considered below (\S~2). Before those propositions can be formulated, however,
the character of the problems posed to quantum mechanics must first be made
clear.

\SourcePage{12}{15}
For this purpose, let us first consider the special relation between quantum
and classical mechanics.

Ordinarily, a more general theory can be formulated as a logically closed
system independently of the less general theory that arises as its limiting
case. Relativistic mechanics, for example, can be constructed from its own
fundamental principles without any reference to Newtonian mechanics. The
fundamental propositions of quantum mechanics, on the other hand, cannot in
principle be formulated without invoking classical mechanics.

The absence of a definite trajectory for an electron\footnote{In this and the
next section, for brevity, we speak of an electron while meaning any quantum
object in general: a particle or system of particles which obeys quantum but
not classical mechanics.} also deprives the electron, considered by itself,
of any other dynamical characteristics\footnote{Here we mean quantities
characterizing the electron's motion, not those characterizing the electron
as a particle (charge and mass), which are parameters.}. It is therefore clear
that no logically closed mechanics could be constructed for a system
consisting solely of quantum objects. A quantitative description of electron
motion requires the presence of physical objects which obey classical
mechanics to sufficient accuracy. When an electron interacts with a
``classical object,'' the state of that object is, in general, altered. Both the character
and the magnitude of this alteration are governed by the electron's state and can therefore
serve as its quantitative characteristic.

In this connection, the ``classical object'' is usually called an
``apparatus,'' and its interaction with the electron is spoken of as a
``measurement.'' It must be emphasized, however, that this does not mean a
process of ``measurement'' involving a physicist as observer. In quantum
mechanics, measurement means any interaction between a classical and a
quantum object, occurring apart from and independently of any observer. The
recognition of the profound role of the measurement concept in quantum
mechanics is due to \emph{Bohr} (\emph{N.~Bohr}).

We have defined an apparatus as a physical object which obeys classical
mechanics to sufficient accuracy; a body of sufficiently great mass is one
example. It should not be supposed, however, that an apparatus must be
macroscopic. Under suitable conditions an object known to be microscopic can
also play the part of an apparatus, since%
\SourcePage{13}{16}
the meaning of ``with sufficient accuracy'' depends on the specific problem
under consideration. In a Wilson chamber, for instance, one tracks the
electron by its mist trail, whose thickness far exceeds atomic dimensions;
at such a level of accuracy in locating the trajectory, the electron is a
fully classical object.

Quantum mechanics thus occupies a very distinctive place among physical
theories: it contains classical mechanics as its limiting case and, at the
same time, requires that limiting case for its own foundation.

We can now formulate the problem addressed by quantum mechanics. A typical
problem is to predict the outcome of a repeated measurement from the known
outcomes of earlier measurements. We shall also see below that, compared with
classical mechanics, quantum mechanics generally restricts the set of values
that various physical quantities (energy, for example) may assume, that is,
the values that may be found by measuring the quantity in question. The
formalism of quantum mechanics must make it possible to determine these
allowed values.

In quantum mechanics the measurement process has an essential peculiarity: it
always acts upon the electron being measured, and for a given measurement
accuracy this action cannot in principle be made arbitrarily weak. The more
accurate the measurement, the stronger its action; only in measurements of
very low accuracy can the action on the measured object be weak. This property
of measurements is logically connected with the fact that the dynamical
characteristics of an electron arise only through the measurement itself. It
is clear that, if the action of the measurement process on the object could be
made arbitrarily weak, the quantity being measured would have a definite value
by itself, irrespective of the measurement.

Of the various types of measurement, the principal one is the measurement of
an electron's coordinates. Within the domain of applicability of
quantum mechanics, the electron's coordinates can always be
measured\footnote{We emphasize once again that by a ``measurement made'' we
mean the interaction of the electron with a classical ``apparatus''; this in
no way presupposes the presence of an outside observer.} to any accuracy.

Suppose that successive measurements of the coordinates%
\SourcePage{14}{17}
of an electron are made at specified time intervals $\Delta t$. As a rule,
these results will fail to follow any smooth curve. Quite the opposite:
the greater the measurement precision, the more abrupt and erratic
the succession of results becomes, consistent with the absence of a path
concept for the electron. A more or less smooth path is obtained only when the
electron's coordinates are measured with a low degree of accuracy, for
example, by the condensation of droplets of vapour in a Wilson chamber.

If instead, while keeping the measurement accuracy unchanged, the intervals
$\Delta t$ between measurements are reduced, adjacent measurements will, of
course, give nearby coordinate values. Yet although the results of a series of
successive measurements will lie in a small region of space, within that
region they will be distributed quite irregularly and will by no means fit any
smooth curve. In particular, as $\Delta t$ tends to zero, the results of
nearby measurements do not tend at all to lie on one straight line.

This last observation makes plain that a classical particle velocity does not
exist in quantum mechanics: the ratio of the coordinate change between two
time instants to the separation $\Delta t$ tends to no well-defined limit. As we shall see later, however, quantum mechanics
does admit a meaningful definition of a particle's velocity at a given
instant, one that reduces to the classical velocity upon passing to classical
mechanics.

In classical mechanics a particle possesses, at each instant, definite
coordinates and velocity; in quantum mechanics the situation is entirely
different. If a measurement has given an electron definite coordinates, it
then has no definite velocity at all. Conversely, an electron having a
definite velocity cannot have a definite position in space. Indeed, the
simultaneous existence of coordinates and velocity at any instant would imply
the presence of a definite velocity, which the electron does not have. Thus,
in quantum mechanics the electron's coordinates and velocity are quantities
that cannot be measured exactly at the same time, that is, cannot
simultaneously have definite values. One may say that the electron's
coordinates and velocity are quantities which do not exist simultaneously. A
quantitative relation governing the possibility of measuring coordinates and
velocity inexactly at the same instant will be derived below.
\SourcePage{15}{18}
In classical mechanics, the state of a physical system is completely described
by specifying all its coordinates and velocities at a given instant; from
these initial data the equations of motion fully determine the system's
behaviour at all future instants. In quantum mechanics
such a description is impossible in principle, since the coordinates and
their corresponding velocities do not exist simultaneously. Hence the state
of a quantum system is described by fewer quantities than in classical
mechanics, that is, its description is less detailed than a classical one.

From this follows a very important consequence regarding the character of
predictions in quantum mechanics. Whereas a classical description is
sufficient to predict with complete precision the future motion of a
mechanical system, the less detailed quantum-mechanical description manifestly
cannot be adequate for this purpose.
Consequently, even when an electron is in a state described as completely as
possible, its behaviour at subsequent instants is in principle not unique.
Quantum mechanics therefore cannot make strictly definite predictions about
the electron's future behaviour. For a specified initial state of an electron,
a subsequent measurement may give different results. The task of quantum
mechanics is only to determine the probability that one or another result will
be obtained in that measurement. In some cases, of course, the probability of
a particular measurement result may equal unity, that is, become a certainty,
so that the result of the measurement is unique.

Every measurement process in quantum mechanics falls into one of two
categories. The first, which includes the majority of measurements, comprises
those that yield no determinate result with certainty for any state of the
system. The second comprises measurements such that for every possible result
there exists a state in which that result is obtained with certainty. It is
precisely these latter measurements, which may be called
\emph{predictable}, that have the principal role in quantum mechanics. The
quantitative characteristics of a state determined by such measurements are
what quantum mechanics calls physical quantities. If in some state a
measurement gives a unique result with certainty, we shall say that the
corresponding physical quantity has a definite value in that state. Below we
shall everywhere%
\SourcePage{16}{19}
understand the expression ``physical quantity'' in precisely the sense stated
here.

We shall repeatedly find below that by no means every collection of physical
quantities in quantum mechanics can be measured simultaneously, that is, can
have definite values simultaneously (we have already discussed one example:
the velocity and coordinates of an electron).

Sets of physical quantities with the following property have an important
role in quantum mechanics: the quantities can be measured simultaneously and,
when they simultaneously have definite values, no other physical quantity
(unless it is a function of them) can have a definite value in that state. We
shall refer to such sets of physical quantities as \emph{complete sets}.

Every description of an electron's state arises from some measurement. We now
state what a complete description of a state means in quantum mechanics.
Completely described states arise from the simultaneous measurement of a
complete set of physical quantities. From the results of such a measurement
one can, in particular, determine the probabilities of the results of any
subsequent measurement, independently of everything that happened to the
electron before the first measurement.

Below, everywhere except in \secref{14}, by states of a quantum system we shall mean
states described in precisely this complete manner.

\section{The superposition principle}

\SourcePage{16}{19}
The radical change in physical ideas about motion that quantum mechanics
requires in comparison with classical mechanics naturally calls for an equally
radical change in the theory's mathematical formalism. In this connection the
first question that arises concerns the means of describing a state in quantum
mechanics.

Let $q$ stand for the totality of coordinates of a quantum
system, and $dq$ for the product of their differentials (it
is called a volume element of the system's \textit{configuration space}); for
one particle, $dq$ coincides with the volume element $dV$ of ordinary space.

The mathematical formalism of quantum mechanics rests on the assertion that a
system's state may be described by a certain (generally complex) function of
the coordinates $\Psi(q)$, and that the squared modulus of this function
determines the probability distribution of the coordinate values:
$|\Psi|^2\,dq$ is%
\SourcePage{17}{20}
the probability that a measurement made on the system will find coordinate
values in the element $dq$ of configuration space. The function $\Psi$ is
called the system's \textit{wave function}\footnote{It was first introduced
into quantum mechanics by \textit{Schrödinger} (\textit{E.~Schrödinger},
1926).}.

Knowledge of the wave function makes it possible in principle to calculate
the probabilities of the various outcomes of any measurement in general (not
necessarily a coordinate measurement). Every one of these probabilities is
given by an expression bilinear in $\Psi$ and $\Psi^*$. In full generality,
such an expression has the form
\begin{equation}
  \iint \Psi(q)\Psi^*(q')\varphi(q,q')\,dq\,dq',
  \tag{2.1}
\end{equation}
where the function $\varphi(q,q')$ depends on the kind and outcome of the
measurement, and the integrations extend over the entire configuration space.
The probability $\Psi\Psi^*$ itself for the different coordinate values is
also an expression of this type\footnote{It is obtained from (2.1) by taking
$\varphi(q,q')=\delta(q-q_0)\delta(q'-q_0)$, where $\delta$ denotes the
so-called $\delta$-function defined below, in \secref{5}; $q_0$ denotes the
coordinate value whose probability is sought.}.

As time passes, the system's state, and with it the wave function, generally
changes. In this sense the wave function may also be regarded as a function of
time. If the wave function is known at some initial instant, then, by the very
meaning of a complete description of a state, it is thereby determined in
principle at every future instant. The actual time dependence of the wave
function is specified by equations to be derived below.

By definition, the sum of the probabilities of all possible values of the
system's coordinates must equal unity. The integral of $|\Psi|^2$ over the
whole configuration space must therefore equal unity:
\begin{equation}
  \int |\Psi|^2\,dq=1.
  \tag{2.2}
\end{equation}
This equality constitutes what is known as the \textit{normalization condition}
on wave functions. When the integral of $|\Psi|^2$ is convergent, by choosing a
suitable constant coefficient, the function $\Psi$ can always be normalized,
as it is said. We shall see below, however, that the integral of $|\Psi|^2$
may diverge, in which case $\Psi$ cannot be normalized by condition (2.2).
\SourcePage{18}{21}
Here $|\Psi|^2$ naturally does not yield the absolute coordinate
probabilities; the ratio of $|\Psi|^2$ between two distinct points of
configuration space, however, does determine the relative probability of
the coordinate values.

Since all quantities with direct physical meaning that are calculated using
the wave function have the form (2.1), in which $\Psi$ occurs multiplied by
$\Psi^*$, it is clear that a normalized wave function is defined only up to a
constant \textit{phase factor} of the form $e^{i\alpha}$, where $\alpha$ is
any real number. This ambiguity is fundamental and cannot be removed; it is,
however, immaterial, since it has no effect on any physical result.

The positive content of quantum mechanics is founded on a number of assertions
about the properties of the wave function, as follows.

Suppose that in the state with wave function $\Psi_1(q)$ a certain measurement
leads with certainty to a definite outcome, outcome 1, while in the state
$\Psi_2(q)$ it leads to outcome 2. It is then postulated that every linear
combination of $\Psi_1$ and $\Psi_2$, that is, every function of the form
$c_1\Psi_1+c_2\Psi_2$ ($c_1$, $c_2$ being constants), describes a state in
which the same measurement yields either outcome 1 or outcome 2. Moreover, if
the time dependence of the states is known, being given in one case by the
function $\Psi_1(q,t)$ and in the other by $\Psi_2(q,t)$, then any linear
combination of them also gives a possible time dependence of a state.

These assertions constitute the so-called \textit{superposition principle for
states}, the fundamental positive principle of quantum mechanics. It follows
from this principle, in particular, that every equation satisfied by wave
functions must be linear in $\Psi$.

Consider a system made up of two parts, and suppose its state is specified in
such a way that each part is described completely\footnote{This, of course,
also gives a complete description of the state of the system as a whole. We
emphasize, however, that the converse statement is by no means valid: in
general, a complete description of the state of the system as a whole does not
yet describe the states of its separate parts completely (see also \secref{14}).}.
One can then assert that the coordinate probabilities for $q_1$ in the
first part are independent of those for the coordinates $q_2$
of the second, and therefore that the probability distribution for the whole
system must equal the product of the probability distributions for its parts.
It follows that the system's wave function $\Psi_{12}(q_1,q_2)$ may be written
as a product of the wave functions $\Psi_1(q_1)$ and
$\Psi_2(q_2)$ of its%
\SourcePage{19}{22}
parts:
\begin{equation}
  \Psi_{12}(q_1,q_2)=\Psi_1(q_1)\Psi_2(q_2).
  \tag{2.3}
\end{equation}
If the two parts do not interact with each other, this relation between the
wave functions of the system and of its parts is preserved at future instants
as well:
\begin{equation}
  \Psi_{12}(q_1,q_2,t)=\Psi_1(q_1,t)\Psi_2(q_2,t).
  \tag{2.4}
\end{equation}

\SourcePage{19}{22}
\section{Operators}

Consider some physical quantity $f$ that characterizes the state of a quantum
system. Strictly speaking, in what follows we ought to speak not of a single
quantity but of an entire complete set of them at once. However, none of the
reasoning is thereby essentially altered, and for the sake of brevity and
simplicity we speak below of only a single physical quantity.

Those values which a given physical quantity may assume are termed its
\emph{eigenvalues} in quantum mechanics, and their aggregate is referred to as
the eigenvalue \emph{spectrum} of that quantity. In classical mechanics,
quantities range, in general, over a continuous series of values. In quantum
mechanics too there exist physical quantities (for example, coordinates) whose
eigenvalues fill a continuous range; in such cases one speaks of a
\emph{continuous spectrum} of eigenvalues. Alongside these quantities, however,
quantum mechanics also possesses quantities whose eigenvalues constitute a
discrete set; one then speaks of a \emph{discrete spectrum}.

For simplicity we shall first assume that the quantity $f$ under consideration
has a discrete spectrum; the case of a continuous spectrum is considered in
\secref{5}. We denote the eigenvalues of $f$ by $f_n$, where the index $n$ runs
through the values $0, 1, 2, 3, \ldots$ We denote the wave function of the
system in the state in which $f$ has the value $f_n$ by $\Psi_n$. The wave
functions $\Psi_n$ bear the name \emph{eigenfunctions} of the physical
quantity $f$. Every such function is taken to be normalized:
\begin{equation}
  \int |\Psi_n|^2\,dq = 1.
  \tag{3.1}
\end{equation}

If the system is in some arbitrary state with wave function $\Psi$, then a
measurement of $f$ performed on it will yield as a result one of the
eigenvalues $f_n$.%
\SourcePage{20}{23}
By the superposition principle, the wave
function $\Psi$ must be a linear combination of those eigenfunctions $\Psi_n$
that correspond to the values $f_n$ which can be found with non-zero probability
in a measurement performed on the system in the state under consideration.
In the general case of an arbitrary state, $\Psi$ may accordingly be written as the series
\begin{equation}
  \Psi = \sum_n a_n\Psi_n,
  \tag{3.2}
\end{equation}
where the summation is over all $n$, and $a_n$ are certain constant
coefficients.

Thus we arrive at the conclusion that every wave function can be, as one says,
expanded in the eigenfunctions of any physical quantity. A system of functions
in terms of which such an expansion can be carried out is spoken of as a
\emph{complete system of functions}.

The expansion (3.2) makes it possible to determine the probabilities of
finding (by means of measurements) one or another value $f_n$ of $f$ for a
system in the state with wave function $\Psi$. Indeed, according to what was
said in the preceding section, these probabilities must be determined by certain
expressions bilinear in $\Psi$ and $\Psi^*$ and must therefore be bilinear in
$a_n$ and $a_n^*$. Moreover, positivity of these expressions is obviously
required. As for the probability associated with a given $f_n$: it must
become unity for a system in the state $\Psi=\Psi_n$ and must vanish whenever
expansion (3.2) of $\Psi$ lacks a term containing the corresponding $\Psi_n$.
The only essentially positive quantity satisfying this condition is the squared
modulus of the coefficient $a_n$. Thus we arrive at the result that the squared
modulus $|a_n|^2$ of every expansion coefficient in (3.2) gives the
probability of the respective value $f_n$ of $f$ in the state with wave
function $\Psi$. The sum of the probabilities of all possible values $f_n$ must
equal unity; in other words, the relation
\begin{equation}
  \sum_n |a_n|^2 = 1.
  \tag{3.3}
\end{equation}
must hold.

If $\Psi$ were not normalized, then the relation (3.3) would also not hold. The
sum $\sum_n |a_n|^2$ would then have to be determined by some expression
bilinear in $\Psi$ and $\Psi^*$ and reducing to unity for normalized $\Psi$.
\SourcePage{21}{24}
The only such expression is the integral $\int\Psi\Psi^*\,dq$. Thus the
equality
\begin{equation}
  \sum_n a_n a_n^* = \int \Psi\Psi^*\,dq.
  \tag{3.4}
\end{equation}
must hold.

On the other hand, multiplying the expansion
$\Psi^*=\sum_n a_n^*\Psi_n^*$ of the complex conjugate function $\Psi^*$ by
$\Psi$ and integrating, we obtain
\[
  \int \Psi\Psi^*\,dq
  = \sum_n a_n^*\int \Psi_n^*\Psi\,dq.
\]
Comparing this with (3.4), we have
\[
  \sum_n a_n a_n^*
  = \sum_n a_n^*\int \Psi_n^*\Psi\,dq,
\]
from which we find the following formula determining the coefficients $a_n$ of
the expansion of $\Psi$ in the eigenfunctions $\Psi_n$:
\begin{equation}
  a_n = \int \Psi\Psi_n^*\,dq.
  \tag{3.5}
\end{equation}

Substituting (3.2) here, we obtain
\[
  a_n = \sum_m a_m\int \Psi_m\Psi_n^*\,dq,
\]
whence it is seen that the eigenfunctions must satisfy the conditions
\begin{equation}
  \int \Psi_m\Psi_n^*\,dq = \delta_{nm},
  \tag{3.6}
\end{equation}
where $\delta_{nm}=1$ for $n=m$ and $\delta_{nm}=0$ for $n\ne m$. The fact that
the integrals of the products $\Psi_m\Psi_n^*$ with $m\ne n$ vanish is spoken of
as the mutual \emph{orthogonality} of the functions $\Psi_n$. The collection of
eigenfunctions $\Psi_n$ thus constitutes a complete system of normalized and mutually
orthogonal (or, as one says for short,~\emph{orthonormal}) functions.

We introduce the concept of the \emph{mean value} $\bar f$ of a quantity $f$ in
a given state. Following the standard definition,
$\bar f$ is the sum of every eigenvalue $f_n$ of the quantity in question, weighted by the respective%
\SourcePage{22}{25}
probability $|a_n|^2$:
\begin{equation}
  \bar f = \sum_n f_n|a_n|^2.
  \tag{3.7}
\end{equation}

We write $\bar f$ in the form of an expression that involves not the expansion
coefficients of $\Psi$ but the function itself. Since (3.7) contains the
products $a_n a_n^*$, it is clear that the desired expression must be bilinear
in $\Psi$ and $\Psi^*$. We introduce a certain mathematical \emph{operator},
which we denote by $\hat f$\footnote{We shall everywhere denote operators by
letters with a hat.}, and define as follows. Let $(\hat f\Psi)$ denote the
result of the action of the operator $\hat f$ on the function $\Psi$. We define
$\hat f$ by requiring that the integral of $(\hat f\Psi)$ multiplied by the
conjugate $\Psi^*$ yield $\bar f$:
\begin{equation}
  \bar f = \int \Psi^*(\hat f\Psi)\,dq.
  \tag{3.8}
\end{equation}

It is easy to see that in the general case the operator $\hat f$ is some linear
integral operator.\footnote{A linear operator is one possessing the properties:
$\hat f(\Psi_1+\Psi_2)=\hat f\Psi_1+\hat f\Psi_2$,
$\hat f(a\Psi)=a\hat f\Psi$, where $\Psi_1$, $\Psi_2$ are arbitrary functions
and $a$ is an arbitrary constant.} Indeed, using the expression (3.5) for
$a_n$, the definition (3.7) of the mean value can be cast as
\[
  \bar f = \sum_n f_n a_n a_n^*
  = \int \Psi^*\left(\sum_n a_n f_n\Psi_n\right)dq.
\]
Comparing with (3.8), we see that the result of the action of the operator
$\hat f$ on the function $\Psi$ has the form
\begin{equation}
  (\hat f\Psi)=\sum_n a_n f_n\Psi_n.
  \tag{3.9}
\end{equation}
Substituting here the expression (3.5) for $a_n$, we find that $\hat f$ is an
integral operator of the form
\begin{equation}
  (\hat f\Psi)=\int K(q,q')\Psi(q')\,dq',
  \tag{3.10}
\end{equation}
where the function $K(q,q')$ (the so-called kernel of the operator) is
\begin{equation}
  K(q,q')=\sum_n f_n\Psi_n^*(q')\Psi_n(q).
  \tag{3.11}
\end{equation}
\SourcePage{23}{26}
Thus, to each physical quantity in quantum mechanics there is assigned a
definite linear operator.

From (3.9) it is seen that if $\Psi$ is one of the eigenfunctions $\Psi_n$ (so
that all $a_n$ except one are zero), then upon the action of the operator
$\hat f$ on it, this function is simply multiplied by the eigenvalue
$f_n$\footnote{Below, wherever it cannot lead to misunderstanding, we shall
omit the parentheses in the expression $(\hat f\Psi)$, the operator being
understood to act on the expression written after it.}:
\begin{equation}
  \hat f\Psi_n=f_n\Psi_n.
  \tag{3.12}
\end{equation}

The eigenfunctions of a physical quantity $f$ are therefore solutions of the
equation
\[
  \hat f\Psi=f\Psi,
\]
in this equation $f$ denotes a constant; the eigenvalues are those particular
values of $f$ that yield solutions obeying the necessary conditions. As will be seen later, the explicit form of operators for the various
physical quantities can be established on direct physical grounds, whereupon
this operator property enables one to obtain eigenfunctions and
eigenvalues by solving the equations $\hat f\Psi=f\Psi$.

For a real physical quantity, both the eigenvalues themselves and their mean
values in any state are real. This fact places a definite constraint on the
properties of the associated operators. Setting the expression (3.8) equal to
its own complex conjugate gives
\begin{equation}
  \int \Psi^*(\hat f\Psi)\,dq
  = \int \Psi(\hat f^*\Psi^*)\,dq,
  \tag{3.13}
\end{equation}
where $\hat f^*$ denotes the operator complex conjugate to
$\hat f$\footnote{By definition, if for the operator $\hat f$ we have
$\hat f\psi=\varphi$, then the complex-conjugate operator $\hat f^*$ is the
operator for which $\hat f^*\psi^*=\varphi^*$.}. For an arbitrary linear
operator such a relation does not, in general, hold, so that it represents a
certain restriction imposed on the possible form of the operators $\hat f$. For
an arbitrary operator $\hat f$ one can define, as one says, the
\emph{transposed} operator $\tilde{\hat f}$, defined%
\SourcePage{24}{27}
so that
\begin{equation}
  \int \Phi(\hat f\Psi)\,dq
  = \int \Psi(\tilde{\hat f}\Phi)\,dq,
  \tag{3.14}
\end{equation}
where $\Psi$, $\Phi$ are two different functions. Choosing as
$\Phi$ the conjugate function $\Psi^*$, comparison with (3.13)
reveals that
\begin{equation}
  \tilde{\hat f}=\hat f^*.
  \tag{3.15}
\end{equation}
An operator that satisfies this condition is called \emph{Hermitian}\footnote{For a
linear integral operator of the form (3.10), Hermiticity requires
that the kernel of the operator satisfy
$K(q,q')=K^*(q',q)$.}. Operators representing real physical quantities within
the mathematical apparatus of quantum mechanics must therefore be Hermitian.

One may also formally consider complex physical quantities,
i.e. quantities possessing complex eigenvalues. Suppose $f$ is one such
quantity; one may then form the complex-conjugate quantity $f^*$, whose eigenvalues are the
complex conjugates of the eigenvalues of $f$. The operator corresponding to the
quantity $f^*$ will be denoted by $\hat f^+$. It is called the \emph{adjoint}
operator to $\hat f$ and must, in general, be distinguished from the
complex-conjugate operator $\hat f^*$. Indeed, by definition of the operator
$\hat f^+$, the mean value of $f^*$ in some state $\Psi$ is
\[
  \overline{f^*}=\int \Psi^*\hat f^+\Psi\,dq.
\]
On the other hand, we have
\[
  {\left(\bar f\right)}^*
  = {\left[\int \Psi^*\hat f\Psi\,dq\right]}^*
  = \int \Psi\hat f^*\Psi^*\,dq
  = \int \Psi^*\tilde{\hat f}^{\,*}\Psi\,dq.
\]
Equating the two expressions, we find that
\begin{equation}
  \hat f^+=\tilde{\hat f}^{\,*},
  \tag{3.16}
\end{equation}
from which it is clear that $\hat f^+$ does not, in general, coincide with
$\hat f^*$. The condition (3.15) can now be written in the form
\begin{equation}
  \hat f=\hat f^+,
  \tag{3.17}
\end{equation}
\SourcePage{25}{28}
i.e. the operator of a real physical quantity coincides with its adjoint
(Hermitian operators are also called \emph{self-adjoint}).

Let us show how one can directly prove the mutual orthogonality of the
eigenfunctions of a Hermitian operator corresponding to different eigenvalues.
Let $f_n$, $f_m$ be two different eigenvalues of the real quantity $f$, and
$\Psi_n$, $\Psi_m$ the corresponding eigenfunctions:
\[
  \hat f\Psi_n=f_n\Psi_n,
  \qquad
  \hat f\Psi_m=f_m\Psi_m.
\]
Multiplying both sides of the first of these equalities by $\Psi_m^*$, and the
complex conjugate of the second equality by $\Psi_n$, and subtracting these
products termwise from each other, we obtain
\[
  \Psi_m^*\hat f\Psi_n-\Psi_n\hat f^*\Psi_m^*
  =(f_n-f_m)\Psi_n\Psi_m^*.
\]
Integrate both sides of this equality over $dq$. Since
$\hat f^*=\tilde{\hat f}$, by virtue of (3.14) the integral of the left-hand
side of the equality vanishes, so that we obtain
\[
  (f_n-f_m)\int\Psi_n\Psi_m^*\,dq=0,
\]
whence, in view of $f_n\ne f_m$, the required orthogonality property of the
functions $\Psi_n$ and $\Psi_m$ follows.

We have been speaking here all along of only a single physical quantity $f$,
whereas we ought to speak, as was noted at the beginning of the section, of a
complete system of simultaneously measurable physical quantities. Then we would
find that to each of these quantities $f$, $g$, $\ldots$ there corresponds its
own operator $\hat f$, $\hat g$, $\ldots$ The eigenfunctions $\Psi_n$
correspond to states in which all the quantities under consideration have
definite values, i.e. correspond to definite sets of eigenvalues
$f_n$, $g_n$, $\ldots$, and are joint solutions of the system of equations
\[
  \hat f\Psi=f\Psi,
  \qquad
  \hat g\Psi=g\Psi,
  \ldots
\]

\section{Addition and Multiplication of Operators}

\SourcePage{25}{28}
If $\hat f$ and $\hat g$ are the operators corresponding to two physical
quantities $f$ and $g$, then $\hat f+\hat g$ corresponds to their sum $f+g$.
In quantum mechanics, however, the meaning of adding different physical
quantities differs essentially according as the quantities are or are not
measurable simultaneously. If $f$ and $g$
\SourcePage{26}{29}
are simultaneously measurable, the operators $\hat f$ and $\hat g$ have
common eigenfunctions. These are also eigenfunctions of $\hat f+\hat g$, whose
eigenvalues are the sums $f_n+g_n$.

If $f$ and $g$ cannot have definite values simultaneously, their sum $f+g$
has a more restricted meaning. All that can be stated is that its mean value
in an arbitrary state equals the sum of the separate mean values of the two
terms:
\begin{equation}
  \overline{f+g}=\bar f+\bar g.
  \tag{4.1}
\end{equation}
The eigenvalues and eigenfunctions of $\hat f+\hat g$, on the other hand, will
in general have no relation to those of $f$ and $g$. Clearly, when $\hat f$
and $\hat g$ are Hermitian, so is $\hat f+\hat g$; its eigenvalues are
therefore real and are the eigenvalues of the new quantity $f+g$ defined in
this way.

We note the following theorem. Let $f_0$ and $g_0$ be the least eigenvalues of
$f$ and $g$, and let $(f+g)_0$ denote the corresponding quantity for $f+g$.
Then
\begin{equation}
  (f+g)_0\geq f_0+g_0.
  \tag{4.2}
\end{equation}
The equality sign applies if $f$ and $g$ are simultaneously measurable. The
proof follows from the evident fact that a quantity's mean value is always at
least its least eigenvalue. In the state in which $(f+g)$ has the value
$(f+g)_0$, we have $\overline{(f+g)}=(f+g)_0$; on the other hand,
$\overline{(f+g)}=\bar f+\bar g\geq f_0+g_0$, which gives (4.2).

Suppose again that $f$ and $g$ are simultaneously measurable. Besides their
sum, we may introduce their product as a quantity whose eigenvalues are the
products of the eigenvalues of $f$ and $g$. The corresponding operator is
readily seen to act by applying one operator to a function and then the other.
Mathematically it is written as the product of $\hat f$ and $\hat g$. Indeed,
if $\Psi_n$ are common eigenfunctions of $\hat f$ and $\hat g$, then
\[
  \hat f\hat g\Psi_n=\hat f(\hat g\Psi_n)
  =\hat f g_n\Psi_n=g_n\hat f\Psi_n=g_nf_n\Psi_n
\]
\SourcePage{27}{30}
(the symbol $\hat f\hat g$ denotes the operator which acts on a function
$\Psi$ first by applying $\hat g$ to $\Psi$, and then by applying $\hat f$ to
the function $\hat g\Psi$). We could equally well have used $\hat g\hat f$,
whose factors occur in the opposite order. Plainly, the actions of these two
operators on the functions $\Psi_n$ give the same result. Since every wave
function $\Psi$ is expandable in the functions $\Psi_n$, it follows that
$\hat f\hat g$ and $\hat g\hat f$ yield the same result when acting on
an arbitrary function. This fact may be written symbolically as
$\hat f\hat g=\hat g\hat f$, or
\begin{equation}
  \hat f\hat g-\hat g\hat f=0.
  \tag{4.3}
\end{equation}

Two such operators $\hat f$ and $\hat g$ are said to \emph{commute} with one
another. We have therefore reached an important result: if two quantities $f$
and $g$ can have definite values simultaneously, their operators commute.

The converse theorem can also be proved (see \secref{11}): when $\hat f$ and
$\hat g$ commute, all their eigenfunctions can be chosen in common, which
physically means that the corresponding physical quantities are simultaneously
measurable. Thus commutation of the operators is a necessary and sufficient
condition for simultaneous measurability of physical quantities.

Raising an operator to a power is a special case of multiplying operators. It
follows from what has been said that the eigenvalues of $\hat f^p$ (where $p$
is an integer) are the eigenvalues of $\hat f$ raised to the same $p$th power.
More generally, any function of an operator, $\varphi(\hat f)$, may be defined
as the operator whose eigenvalues are the same function $\varphi(f)$ of the
eigenvalues of $\hat f$. When $\varphi(f)$ has a Taylor-series expansion, this
expansion reduces the action of $\varphi(\hat f)$ to the actions of the
various powers $\hat f^p$.

In particular, $\hat f^{-1}$ is called the operator \emph{inverse} to
$\hat f$. Clearly, successive application of $\hat f$ and $\hat f^{-1}$ to an
arbitrary function leaves it unchanged; that is,
$\hat f\hat f^{-1}=\hat f^{-1}\hat f=1$.
\SourcePage{28}{31}
When $f$ and $g$ are not simultaneously measurable, their product has no such
direct meaning. This is already apparent because $\hat f\hat g$ is then not
Hermitian and consequently cannot correspond to a real physical quantity.
Indeed, by the definition of the transposed operator,
\[
  \int \Psi\hat f\hat g\Phi\,dq
  =\int \Psi\hat f(\hat g\Phi)\,dq
  =\int(\hat g\Phi)(\tilde f\Psi)\,dq.
\]
Here $\tilde f$ acts only on $\Psi$, while $\hat g$ acts on $\Phi$; the
integrand is therefore simply the product of the two functions $\hat g\Phi$
and $\tilde f\Psi$. Applying the definition of the transposed operator
once more gives
\[
  \int \Psi\hat f\hat g\Phi\,dq
  =\int(\tilde f\Psi)(\hat g\Phi)\,dq
  =\int\Phi\tilde g\,\tilde f\Psi\,dq.
\]
We have thus obtained an integral in which $\Psi$ and $\Phi$ have exchanged
places relative to the original one. In other words,
$\tilde g\,\tilde f$ is the operator transposed to $\hat f\hat g$, so
we may write
\begin{equation}
  \tilde{fg}=\tilde g\,\tilde f,
  \tag{4.4}
\end{equation}
that is, the transpose of the product $\hat f\hat g$ is the product of the
transposed factors written in reverse order. Taking the complex conjugate of
both sides of (4.4), we find
\begin{equation}
  (\hat f\hat g)^+=\hat g^+\hat f^+.
  \tag{4.5}
\end{equation}
If both $\hat f$ and $\hat g$ are Hermitian, then
$(\hat f\hat g)^+=\hat g\hat f$. Hence $\hat f\hat g$ is Hermitian only when
the factors $\hat f$ and $\hat g$ commute.

From the products $\hat f\hat g$ and $\hat g\hat f$ of two noncommuting
Hermitian operators one can nevertheless form another Hermitian operator,
their \emph{symmetrized product}
\begin{equation}
  \frac{1}{2}(\hat f\hat g+\hat g\hat f).
  \tag{4.6}
\end{equation}

It is also readily verified that the difference
$\hat f\hat g-\hat g\hat f$ is an ``anti-Hermitian'' operator (that is, its
transpose is equal to the negative of its complex
\SourcePage{29}{32}
conjugate). Multiplication by $i$ makes it Hermitian; hence
\begin{equation}
  i(\hat f\hat g-\hat g\hat f)
  \tag{4.7}
\end{equation}
is likewise a Hermitian operator.

For brevity, later on we shall sometimes use the notation
\begin{equation}
  \{\hat f,\hat g\}=\hat f\hat g-\hat g\hat f
  \tag{4.8}
\end{equation}
for the so-called \emph{commutator} of the operators. It is easy to verify the
relation
\begin{equation}
  \{\hat f\hat g,\hat h\}
  =\{\hat f,\hat h\}\hat g+\hat f\{\hat g,\hat h\}.
  \tag{4.9}
\end{equation}
Notice that $\{\hat f,\hat h\}=0$ and $\{\hat g,\hat h\}=0$ do not in general
imply that $\hat f$ and $\hat g$ commute with each other.

\section{Continuous spectrum}

\SourcePage{29}{32}
All the relations obtained in \secref{3}, \secref{4} for the properties of discrete-spectrum
eigenfunctions extend without difficulty to a continuous spectrum of
eigenvalues.

Let $f$ be a physical quantity possessing a continuous spectrum. We shall
denote its eigenvalues simply by the same letter $f$ without an index; the
corresponding eigenfunctions will be written $\Psi_f$. In the same way that an
arbitrary wave function $\Psi$ admits the series expansion (3.2) in
eigenfunctions of a discrete-spectrum quantity, it can also be expanded --- this
time as an integral --- over the complete system of eigenfunctions of a
continuous-spectrum quantity. Such an expansion takes the form
\begin{equation}
  \Psi(q)=\int a_f\Psi_f(q)\,df,
  \tag{5.1}
\end{equation}
where the integration is taken over the whole domain of values that the
quantity $f$ can assume.

More complicated than in the discrete-spectrum case is the question of
normalizing the continuous-spectrum eigenfunctions. As we shall see below, the
requirement that the integral of the squared modulus equal unity cannot be
satisfied here. Instead, let us set ourselves
the aim of normalizing the functions $\Psi_f$ in such a way that $|a_f|^2\,df$
represents the probability for the physical quantity under consideration, in
the state described by the wave function $\Psi$, to have a value in the given
interval between $f$ and $f+df$.
\SourcePage{30}{33}
Since the sum of the probabilities of all possible
values of $f$ must be equal to unity, we have
\begin{equation}
  \int |a_f|^2\,df=1
  \tag{5.2}
\end{equation}
(analogously to relation (3.3) for the discrete spectrum).

Proceeding exactly as we did in deriving formula (3.5), and using the same
arguments, we write, on one hand,
\[
  \int\Psi\Psi^*\,dq=\int|a_f|^2\,df,
\]
and, on other hand,
\[
  \int\Psi\Psi^*\,dq
  =\iint a_f^*\Psi_f^*\Psi\,df\,dq.
\]
By comparing the two expressions we find the formula that determines the
expansion coefficients,
\begin{equation}
  a_f=\int\Psi(q)\Psi_f^*(q)\,dq,
  \tag{5.3}
\end{equation}
exactly analogous to (3.5).

To derive the normalization condition we now substitute (5.1) into (5.3):
\[
  a_f=\int a_{f'}\left(\int\Psi_{f'}\Psi_f^*\,dq\right)df'.
\]
This relation must hold for arbitrary $a_f$ and must therefore be satisfied
identically. For this to be so it is first of all necessary that the
coefficient of $a_{f'}$ under the integral sign (i.e. the integral
$\int\Psi_{f'}\Psi_f^*\,dq$) vanish for all $f'\ne f$. For $f'=f$ this
coefficient must become infinite (otherwise the integral over $df'$ would
simply be equal to zero). Thus the integral $\int\Psi_{f'}\Psi_f^*\,dq$ is a
function of the difference $f-f'$ which vanishes for non-zero values of the
argument and becomes infinite when the argument is zero. We denote this function
by $\delta(f'-f)$:
\begin{equation}
  \int\Psi_{f'}\Psi_f^*\,dq=\delta(f'-f).
  \tag{5.4}
\end{equation}

The manner in which the function $\delta(f'-f)$ becomes infinite at
$f'-f=0$ is determined by the requirement that there should be
\[
  \int\delta(f'-f)a_{f'}\,df'=a_f.
\]
\SourcePage{31}{34}
Clearly, for this to hold it is necessary that
\[
  \int\delta(f'-f)\,df'=1.
\]
A function defined in this way is called the
\textit{delta-function}\footnote{The delta-function was introduced into
theoretical physics by \textit{Dirac} (\textit{P.~A.~M. Dirac}).}. Let us write
out once more the formulas that define it. We have
\begin{equation}
  \delta(x)=0\quad\text{at }x\ne0,\qquad \delta(0)=\infty,
  \tag{5.5}
\end{equation}
and moreover in such a way that
\begin{equation}
  \int_{-\infty}^{+\infty}\delta(x)\,dx=1.
  \tag{5.6}
\end{equation}
As the limits of integration one may write any other pair between which the
point $x=0$ lies. If $f(x)$ is a function continuous at $x=0$, then
\begin{equation}
  \int_{-\infty}^{+\infty}\delta(x)f(x)\,dx=f(0).
  \tag{5.7}
\end{equation}
In more general form this formula can be written as
\begin{equation}
  \int\delta(x-a)f(x)\,dx=f(a),
  \tag{5.8}
\end{equation}
where the region of integration includes the point $x=a$, and $f(x)$ is
continuous at $x=a$. It is also obvious that the delta-function is even, i.e.
\begin{equation}
  \delta(-x)=\delta(x).
  \tag{5.9}
\end{equation}
Finally, writing
\[
  \int_{-\infty}^{\infty}\delta(\alpha x)\,dx
  =\int_{-\infty}^{\infty}\delta(y)\frac{dy}{|\alpha|}
  =\frac{1}{|\alpha|},
\]
we arrive at the conclusion that
\begin{equation}
  \delta(\alpha x)=\frac{1}{|\alpha|}\delta(x),
  \tag{5.10}
\end{equation}
where $\alpha$ is any constant.
\SourcePage{32}{35}
Formula (5.4) expresses the normalization rule for the eigenfunctions of the
continuous spectrum; it replaces condition (3.6) of the discrete spectrum. We
see that the functions $\Psi_f$ and $\Psi_{f'}$ with $f\ne f'$ remain
orthogonal to each other. The integrals of the squares $|\Psi_f|^2$ of the
functions of the continuous spectrum, however, diverge.

The functions $\Psi_f(q)$ satisfy one further relation, similar to (5.4). To
derive it we substitute (5.3) into (5.1), which gives
\[
  \Psi(q)=\int\Psi(q')
  \left(\int\Psi_f^*(q')\Psi_f(q)\,df\right)dq',
\]
whence we immediately conclude that one must have
\begin{equation}
  \int\Psi_f^*(q')\Psi_f(q)\,df=\delta(q-q').
  \tag{5.11}
\end{equation}
An analogous relation can, of course, be derived for the discrete spectrum,
where it takes the form
\begin{equation}
  \sum_n\Psi_n^*(q')\Psi_n(q)=\delta(q'-q).
  \tag{5.12}
\end{equation}

On comparing the pair of formulas (5.1) and (5.4) with the pair (5.3) and
(5.11), we see that, on the one hand, the functions $\Psi_f(q)$ effect the
expansion of the function $\Psi(q)$ with the expansion coefficients $a_f$, while,
on the other hand, formula (5.3) can be regarded as a completely analogous
expansion of the function $a_f\equiv a(f)$ in the functions $\Psi_f^*(q)$, with
$\Psi(q)$ playing the role of the expansion coefficients. The function $a(f)$,
like $\Psi(q)$, determines the state of the system completely; it is spoken of
as the wave function in the \textit{$f$-representation} (and the function
$\Psi(q)$ as the wave function in the \textit{$q$-representation}). Just as
$|\Psi(q)|^2$ determines the probability for the system to have coordinates in
a given interval $dq$, so $|a(f)|^2$ determines the probability of values of
the quantity $f$ in a given interval $df$. The functions $\Psi_f(q)$, for their
part, are, on the one hand, the eigenfunctions of the quantity $f$ in the
$q$-representation, while their complex conjugates $\Psi_f^*(q)$ serve as
eigenfunctions of the coordinate $q$ in the $f$-representation.

Suppose $\varphi(f)$ is a function of the quantity $f$, with $\varphi$ and
$f$ standing in one-to-one correspondence. Every function $\Psi_f(q)$ may then
equally be viewed as an eigenfunction of the quantity
$\varphi$. In doing so, however, the normalization of these functions must be
changed. Indeed, the eigen-
\SourcePage{33}{36}
functions $\Psi_\varphi(q)$ of the quantity $\varphi$ must be normalized by the
condition
\[
  \int\Psi_{\varphi(f')}\Psi_{\varphi(f)}^*\,dq
  =\delta[\varphi(f')-\varphi(f)],
\]
whereas the functions $\Psi_f$ are normalized by condition (5.4). The argument
of the delta-function vanishes at $f'=f$. For $f'$ close to $f$ we have
\[
  \varphi(f')-\varphi(f)=\frac{d\varphi(f)}{df}(f'-f).
\]
Making use of expression (5.10), one can write\footnote{In general, if
$\varphi(x)$ is a single-valued function (its inverse, however, may be
many-valued), then the formula
\begin{equation}
  \delta[\varphi(x)]
  =\sum_i\frac{1}{|\varphi'(\alpha_i)|}\delta(x-\alpha_i),
  \tag{5.13a}
\end{equation}
holds, where $\alpha_i$ are the roots of the equation $\varphi(x)=0$.}
\begin{equation}
  \delta[\varphi(f')-\varphi(f)]
  =\frac{1}{|d\varphi(f)/df|}\delta(f'-f).
  \tag{5.13}
\end{equation}
Comparison of (5.13) with (5.4) now shows that the functions $\Psi_\varphi$ and
$\Psi_f$ are related to each other by the relation
\begin{equation}
  \Psi_{\varphi(f)}
  =\frac{1}{\sqrt{|d\varphi(f)/df|}}\Psi_f.
  \tag{5.14}
\end{equation}

There exist physical quantities which, over some domain of their values,
possess a discrete spectrum, and over another possess a continuous one. For the
eigenfunctions of such a quantity, of course, all the same relations that were
derived in this and the preceding sections continue to hold. It need only be
noted that the complete system of functions is formed by the collection of the
eigenfunctions of both spectra together. Therefore the expansion of an
arbitrary wave function in the eigenfunctions of such a quantity has the form
\begin{equation}
  \Psi(q)=\sum_n a_n\Psi_n(q)+\int a_f\Psi_f(q)\,df,
  \tag{5.15}
\end{equation}
where the summation runs over the discrete spectrum, and the integration over
the entire continuous spectrum.

An example of a quantity possessing a continuous spectrum is the coordinate
$q$ itself. One readily sees that the associated operator reduces to
multiplication by $q$. Indeed,
\SourcePage{34}{37}
since the square $|\Psi(q)|^2$ governs the probability of different coordinate
values, the mean value of the coordinate is
\[
  \bar q=\int q|\Psi|^2\,dq=\int\Psi^*q\Psi\,dq.
\]
Comparing this expression with the definition of operators according to (3.8),
we see that\footnote{In what follows we shall agree, for simplicity of notation,
to write operators that reduce to multiplication by some quantity everywhere
simply as that quantity itself.}
\begin{equation}
  \hat q=q.
  \tag{5.16}
\end{equation}
The eigenfunctions of this operator must be determined, according to the general
rule, by the equation $q\Psi_{q_0}=q_0\Psi_{q_0}$, where $q_0$ is used
temporarily to denote concrete values of the coordinate, in distinction from
the variable $q$. Since this equality can be satisfied either when
$\Psi_{q_0}=0$ or when $q=q_0$, it is clear that the normalized eigenfunctions
are\footnote{The expansion coefficients of an arbitrary function $\Psi$ in these
eigenfunctions are $a_{q_0}=\int\Psi(q)\delta(q-q_0)\,dq=\Psi(q_0)$. The
probability of values of the coordinate in a given interval $dq_0$ is
$|a_{q_0}|^2\,dq_0=|\Psi(q_0)|^2\,dq_0$, as it should be.}
\begin{equation}
  \Psi_{q_0}=\delta(q-q_0).
  \tag{5.17}
\end{equation}

\SourcePage{34}{37}
\section{The Limiting Transition}

Classical mechanics is contained in quantum mechanics as a limiting case. We
must determine how this limiting transition takes place.

In quantum mechanics a wave function describes the electron, governing the
various possible values of its coordinate; thus far we know
only that this function solves a certain linear partial differential equation.
Classical mechanics, by contrast, treats the electron as a material particle
moving along a trajectory fixed completely by the equations of motion. A
relation analogous in a certain sense to that between quantum and classical
mechanics occurs in electrodynamics between wave optics and geometrical
optics. Wave optics describes electromagnetic waves by electric- and
magnetic-field vectors satisfying a particular system of linear differential
equations (Maxwell's equations). Geometrical optics, however,%
\SourcePage{35}{38}
describes the propagation of light along definite paths---rays. This analogy
leads one to conclude that the quantum-to-classical limiting transition is
analogous to that from wave optics to geometrical optics.

We now recall how this latter transition is carried out mathematically (see II,
\secref{53}). Let $u$ be any component of the field in an electromagnetic wave. It
may be written $u=ae^{i\varphi}$, with real amplitude $a$ and phase $\varphi$
(the phase is called the eikonal in geometrical optics). The geometrical-optics
limit corresponds to short wavelengths; mathematically, this means that
$\varphi$ changes greatly over short distances. In particular, its absolute
magnitude may be regarded as large.

Correspondingly, assume that the classical limiting case in quantum mechanics
is represented by wave functions of the form $\Psi=ae^{i\varphi}$, where $a$
varies slowly and $\varphi$ takes large values. In mechanics, as is known, a
particle trajectory can be found from a variational principle according to
which the so-called action $S$ of the mechanical system must be a minimum (the
principle of least action). In geometrical optics the ray path is instead
fixed by Fermat's principle: the ``optical path length'' of the ray, that is,
the difference between its phases at the end and at the beginning of the path,
must be a minimum.

This analogy allows us to conclude that, in the classical limiting case, the
phase $\varphi$ of the wave function must be proportional to the mechanical
action $S$ of the physical system under consideration: $S=\operatorname{const}
\cdot\varphi$. The proportionality factor is called \emph{Planck's constant}
and is denoted by $\hbar$\footnote{It was introduced into physics by
\emph{Planck} (\emph{M.~Planck}, 1900). The constant $\hbar$ used throughout
this book is, strictly speaking, Planck's constant $h$ divided by $2\pi$
(Dirac's notation).}. Since $\varphi$ is dimensionless, it has the dimensions
of action, and its value is
\[
  \hbar=1{,}055\cdot10^{-27}\,\text{erg}\cdot\text{s}.
\]

Thus the wave function of an ``almost classical'' (or, as it is called,
\emph{quasiclassical}) physical system has the form
\begin{equation}
  \Psi=ae^{iS/\hbar}.
  \tag{6.1}
\end{equation}
\SourcePage{36}{39}
Planck's constant has a fundamental part in every quantum phenomenon. Its
relative magnitude (in comparison with other quantities having the same
dimensions) determines the ``degree of quantumness'' of a given physical
system. Passage from quantum to classical mechanics corresponds to a large
phase and can formally be described by taking the limit $\hbar\to0$ (just as
the passage from wave to geometrical optics corresponds to the zero-wavelength
limit, $\lambda\to0$).

We have established the limiting form of the wave function, but must still ask
how it is related to classical motion along a trajectory. In general, motion
described by a wave function does not at all turn into motion along one fixed
trajectory. Its connection with classical motion is this: if the wave function,
and hence the probability distribution of the coordinates, is specified at an
initial instant, the distribution thereafter ``moves'' as required by the
laws of classical mechanics (see the end of \secref{17} for details).

To obtain motion along a definite trajectory one must begin with a wave
function of a special kind, appreciably different from zero only within a very
small region of space (a so-called \emph{wave packet}); the dimensions of this
region may be made to tend to zero together with $\hbar$. One can then state
that, in the quasiclassical case, the wave packet moves through space along the
particle's classical trajectory.

Finally, in the limit, quantum-mechanical operators must reduce simply to
multiplication by their corresponding physical quantities.

\SourcePage{36}{39}
\section{The Wave Function and Measurements}

Let us return to the measurement process, whose properties were discussed
qualitatively in \secref{1}, and show how those properties are connected with the
mathematical apparatus of quantum mechanics.

Consider a system made up of two parts: a classical apparatus and an electron
(treated as a quantum object). Measurement consists in bringing these two
parts into interaction. The apparatus consequently passes from its initial
state into some other state, and from this change we draw conclusions about
the state of the electron. States of the apparatus are distinguished by the
values of a physical quantity characterizing it (or%
\SourcePage{37}{40}
quantities)---the apparatus ``reading.'' Denote this quantity provisionally by
$g$, and its eigenvalues by $g_n$. Because the apparatus is classical, these
values generally range over a continuum; solely to simplify the notation in
the formulas below, however, we shall suppose that the spectrum is discrete.
The states of the apparatus are described by quasiclassical wave functions,
which we denote by $\Phi_n(\xi)$; the index $n$ corresponds to the reading
$g_n$, and $\xi$ stands collectively for its coordinates.
That the apparatus is classical manifests itself in the circumstance that at
each instant one can assert with certainty that it occupies one of the known states
$\Phi_n$, with some definite value of $g$. Such a statement would of course
be false for a quantum system.

Let $\Phi_0(\xi)$ be the wave function of the apparatus in its initial state
(before measurement), and let $\Psi(q)$ be an arbitrary normalized initial
wave function of the electron ($q$ denotes its coordinates). These functions
describe apparatus and electron independently, so the initial wave function
of the complete system is the product
\begin{equation}
  \Psi(q)\Phi_0(\xi).
  \tag{7.1}
\end{equation}
The apparatus and electron then interact. In principle, the equations of
quantum mechanics allow the time variation of the system's wave function to
be followed. After the measurement it will plainly no longer be a product of
functions of $\xi$ and $q$. Expanding it in the eigenfunctions $\Phi_n$ of the
apparatus (which form a complete system of functions), we obtain a sum of the
form
\begin{equation}
  \sum_n A_n(q)\Phi_n(\xi),
  \tag{7.2}
\end{equation}
where the $A_n(q)$ are certain functions of $q$.

At this point the ``classical nature'' of the apparatus enters, together with
the double role of classical mechanics as both a limiting case and a basis of
quantum mechanics. As already noted, because the apparatus is classical, the
quantity $g$ (its ``reading'') has a definite value at every instant. It
follows that after measurement the actual state of apparatus plus electron is
described not by the whole sum (7.2), but by only the one term corresponding
to the apparatus reading $g_n$:
\begin{equation}
  A_n(q)\Phi_n(\xi).
  \tag{7.3}
\end{equation}
\SourcePage{38}{41}
Consequently, $A_n(q)$ is proportional to the electron's wave function after
measurement. It is not itself that wave function, as is already evident from
the fact that $A_n(q)$ is not normalized. It contains both information about
the properties of the state produced and the probability, determined by the
system's initial state, that the apparatus will give its $n$th reading.

By the linearity of the equations of quantum mechanics, the relation between
$A_n(q)$ and the initial electron wave function $\Psi(q)$ is in general given
by a linear integral operator,
\begin{equation}
  A_n(q)=\int K_n(q,q')\Psi(q')\,dq',
  \tag{7.4}
\end{equation}
whose kernel $K_n(q,q')$ characterizes the measurement process in question.

We suppose that the measurement under consideration yields a complete
description of the electron's state. In other words (see \secref{1}), in the state
produced, the probabilities for all quantities must be independent of the
electron's preceding state (before measurement). Mathematically, this means
that the form of the functions $A_n(q)$ must be fixed by the measurement
process itself and must not depend on the initial electron wave function
$\Psi(q)$.

The $A_n$ must therefore have the form
\begin{equation}
  A_n(q)=a_n\varphi_n(q),
  \tag{7.5}
\end{equation}
where the $\varphi_n$ are certain functions that we shall take to be
normalized, while only the constants $a_n$ depend on the initial state
$\Psi(q)$. In the integral relation (7.4), this corresponds to a kernel
$K_n(q,q')$ which factors into functions of $q$ and $q'$ separately:
\begin{equation}
  K_n(q,q')=\varphi_n(q)\Psi_n^*(q').
  \tag{7.6}
\end{equation}
The linear dependence of the constants $a_n$ on $\Psi(q)$ is then expressed
by formulas of the form
\begin{equation}
  a_n=\int\Psi(q)\Psi_n^*(q)\,dq,
  \tag{7.7}
\end{equation}
where the $\Psi_n(q)$ are certain functions fixed by the measurement process.

The functions $\varphi_n(q)$ are the normalized wave functions of the
electron after measurement. We thus see how the theory's mathematical
formalism represents the possibility of producing, by a measurement, an
electron state described by a definite wave function.
\SourcePage{39}{42}
When a measurement is made on an electron having a specified wave function
$\Psi(q)$, the constants $a_n$ have a simple physical meaning: by the general
rules, $|a_n|^2$ is the probability that the measurement gives the $n$th
result. The probabilities of all results sum to unity:
\begin{equation}
  \sum_n|a_n|^2=1.
  \tag{7.8}
\end{equation}
For an arbitrary normalized function $\Psi(q)$, the validity of (7.7) and
(7.8) is equivalent (compare \secref{3}) to the statement that $\Psi(q)$ can be
expanded in the functions $\Psi_n(q)$. Hence the functions $\Psi_n(q)$ form
a complete normalized and mutually orthogonal set.

If the electron's initial wave function coincides with one of the functions
$\Psi_n(q)$, the corresponding constant $a_n$ is plainly unity and all the
others vanish. Put differently, if the measurement is performed on an electron
in the state $\Psi_n(q)$, it will with certainty give a definite ($n$th) result.

Taken together, these properties of the functions $\Psi_n(q)$ demonstrate
that they are eigenfunctions of a certain physical quantity that characterizes the electron (call it
$f$), and the measurement being considered may be called a measurement of
that quantity.

It is essential that, in general, the functions $\Psi_n(q)$ do not coincide
with the functions $\varphi_n(q)$ (indeed, the latter need not even be
mutually orthogonal or form an eigenfunction system of any operator). This
fact expresses first of all the non-repeatability of measurement results in
quantum mechanics. If the electron was in the state $\Psi_n(q)$, a measurement
of $f$ made upon it will reveal the value $f_n$ with certainty. After the
measurement, however, the electron is in a different state $\varphi_n(q)$, in
which $f$ need no longer have any definite value. Thus, if a second
measurement were made immediately after the first, the value found for $f$
would not coincide with the result of the first
measurement\footnote{There is, however, an important exception to the
non-repeatability of measurements: coordinate is the one quantity whose
measurement is repeatable. Two measurements of an electron's coordinate made
at a sufficiently short interval must give nearby values; otherwise the
electron would have infinite velocity. The mathematical origin of this
exception lies in the commutativity of the coordinate with the
electron--apparatus interaction energy operator, which (in the nonrelativistic
theory) depends solely on
coordinates.}. To predict (in the sense of%
\SourcePage{40}{43}
calculating the probability of) the result of a repeated measurement when the
first result is known, one takes from the first measurement the wave function
$\varphi_n(q)$ of the state it produced, and from the second the wave function
$\Psi_m(q)$ of the state whose probability is sought. More explicitly, the
equations of quantum mechanics determine a wave function
$\varphi_n(q,t)$ which equals $\varphi_n(q)$ at the time of the first
measurement. If the second measurement is made at time $t$, the probability
of its $m$th result is the squared modulus of the integral
$\int\varphi_n(q,t)\Psi_m^*(q)\,dq$.

We see that measurement in quantum mechanics is ``two-faced'': its roles with
respect to past and future do not coincide. Toward the past it ``verifies''
the probabilities of the various possible results predicted from the state
created by the preceding measurement. Toward the future it creates a new
state (see also \secref{44}). A profound irreversibility is therefore inherent in
the measurement process itself.

This irreversibility has important consequences of principle. As will be seen
later (see the end of \secref{18}), the fundamental equations of quantum mechanics
by themselves are symmetric under reversal of the sign of time; in this
respect quantum mechanics does not differ from classical mechanics. The
irreversibility of measurement, however, makes the two directions of time
physically inequivalent in quantum phenomena, and thus produces a distinction
between future and past.

